% Generated draft: requires mathematical and source review.
\documentclass[11pt]{article}
\usepackage[margin=25mm]{geometry}
\usepackage{fontspec}
\setmainfont{Latin Modern Roman}
\usepackage{amsmath,amssymb,mathtools,unicode-math}
\setmathfont{Latin Modern Math}
\usepackage{enumitem,needspace,xurl,hyperref}
\setlength{\parindent}{0pt}
\setlength{\parskip}{5pt}
\setlength{\emergencystretch}{3em}
\newcommand{\N}{\mathbb{N}}
\newcommand{\Z}{\mathbb{Z}}
\newcommand{\Q}{\mathbb{Q}}
\newcommand{\R}{\mathbb{R}}
\newcommand{\C}{\mathbb{C}}
\newcommand{\F}{\mathbb{F}}
\newcommand{\A}{\mathbb{A}}
\newcommand{\PP}{\mathbb{P}}
\newcommand{\E}{\mathbb{E}}
\newcommand{\eps}{\varepsilon}
\newcommand{\dd}{\,\mathrm{d}}
\begin{document}
\section*{UnsolvedMath ID: 30001114}
\textbf{Source problem code:} catalogue entry
\subsection{Definitions and mathematical statement}
Let \(p>0\) be a prime, let \(X\) be a regular scheme of dimension two in characteristic \(p\), and let a finite group \(G\) act on \(X\). Assume that the scheme quotient \(Y=X/G\) exists, and let \(y\in Y\) be an isolated singular point: \(\mathcal O_{Y,y}\) is not regular, while some open neighborhood of \(y\) has no other singular points. The supplied statement imposes \(p\mid |G|\) and describes the singularities under consideration as wild.

\textbf{Requested classification.} Classify all singularities \((Y,y)\) arising under these conditions up to analytic isomorphism. The source also requests a classification of the possible resolution graphs and associated combinatorial data of their minimal normal-crossing resolutions.

Here a resolution of \((Y,y)\), when defined and available in the intended category, is a proper birational morphism \(f:\widetilde U\to U\), where \(U\) is a neighborhood of \(y\), \(\widetilde U\) is regular, and \(f\) is an isomorphism over \(U\setminus\{y\}\). The normal-crossing requirement concerns the reduced exceptional divisor \(E=(f^{-1}(y))_{\mathrm{red}}\). Its resolution graph records the irreducible exceptional components and their intersections; numerical decorations, such as genera and self-intersection numbers, require the relevant geometric hypotheses and conventions.

\textbf{Unresolved qualifications.} The record does not specify a base field, Noetherian or finite-type hypotheses, residue fields, or assumptions ensuring the intended quotient and resolution theory. It also does not define analytic isomorphism in this positive-characteristic scheme setting. One possible formal interpretation would compare completed local rings \(\widehat{\mathcal O}_{Y,y}\), with any required base-field structure preserved, but this interpretation is not selected by the record. The meanings of minimality, normal crossings versus strict normal crossings, and the full associated combinatorial data are likewise unspecified.

The condition \(p\mid |G|\) concerns the whole acting group; it does not by itself assert that a point above \(y\) has stabilizer of order divisible by \(p\). Thus the intended local meaning of wildness requires clarification. Finally, the source calls analytic classification and classification by resolution graphs and combinatorial data equivalent. No completeness theorem or sufficiently detailed invariant is supplied to justify this equivalence. Accordingly, the analytic classification and the classification of resolution data are retained as requested targets, without asserting that the latter determines the former.
\subsection{Short English statement}
Classify the isolated quotient surface singularities described by the source in characteristic \(p>0\), with \(p\mid |G|\), and describe their minimal normal-crossing resolution data. The intended analytic category, local wildness condition, geometric hypotheses, and claimed equivalence with classification by resolution data remain unresolved.
\subsection{Sources}
\begin{itemize}
\item Supplied local UnsolvedMath registry record, permanent ID \(30001114\), source problem code \texttt{catalogue entry}, title \emph{Classification of Wild Quotient Surface Singularities}, in the Oberwolfach Reports Open Problems collection. The supplied record includes original and curated statement fields; its \texttt{corrected\_verified} label is a supplied curation label, not independent mathematical verification.
\item Supplied source citation: \emph{Mini-Workshop: Group Actions on Curves: Reduction and Lifting (2009)}, Oberwolfach Reports, pp.\ 3029--3064, DOI \texttt{10.4171/owr/2008/54}. \url{https://doi.org/10.4171/owr/2008/54}. The background attributes the collection to workshop open-problem discussions. No individual proposer is supplied. The proposed year and year tag are 2008, while the citation labels the workshop 2009; this discrepancy is unresolved.
\item Supplied literature source, labeled ``original OWR report'': \url{https://doi.org/10.4171/owr/2008/54}.
\item Supplied literature source, labeled ``wild quotient surface singularities'': \url{https://arxiv.org/abs/1209.3605}.
\item Supplied literature source, labeled ``wild cyclic quotient singularities'': \url{https://arxiv.org/abs/2101.08090}.
\item Literature assessment dated exactly \texttt{2026-08-21}: the supplied curation labels the problem partially solved and reports substantial classification results for cyclic and surface families, including explicit resolution graphs, while reporting that a classification of all wild singularities remains unavailable. This is supplied dated curation derived from the OWR collection, not an independently verified literature assessment; the status and cited claims require independent verification.
\end{itemize}
\end{document}
